\section*{अध्याय \ref{ch:b3:spin-two-level} --- चक्रण और द्वि-स्तरीय तंत्र}

\begin{solution}{exo:b3:spin-two-level:1}
(क) सीधा गुणन। (ख) $[\hat S_x, \hat S_y] =
(\hbar/2)^2\,2\iu\sigma_z = \iu\hbar\hat S_z$। (ग) $\hat S^2 =
(\hbar/2)^2 \times 3\,\mathbb 1 = \tfrac34\hbar^2$: अर्थात् सचमुच
$\tfrac12\cdot\tfrac32\hbar^2$। (घ) $\{\mathbb 1, \sigma_x,
\sigma_y, \sigma_z\}$ हर्मिटीय $2\times2$ आव्यूहों की चारों वास्तविक
विमाएँ फैलाते हैं ($a + \vect b\cdot\vect\sigma$, जिसमें $a,
\vect b$ वास्तविक हैं)।
\end{solution}

\begin{solution}{exo:b3:spin-two-level:2}
(क) $\ket{\pm x} = (\ket\uparrow \pm \ket\downarrow)/\sqrt2$;
$\ket{\pm y} = (\ket\uparrow \pm \iu\ket\downarrow)/\sqrt2$। (ख)
आधे-आधे। (ग) फिर आधे-आधे: $|\braket{\pm y}{+x}|^2 =
|1 \mp \iu|^2/4 = \tfrac12$। (घ) $\ket{+x}$: विषुवत् पर, $\varphi =
0$ पर; और $z$ या $y$ के अनुदिश मापने पर वह $90^\circ$ दूर के ध्रुवों
पर प्रक्षिप्त होता है --- अतः सदा $\cos^2(45^\circ) = \tfrac12$।
\end{solution}

\begin{solution}{exo:b3:spin-two-level:3}
(क) $\cos^2(\theta/2)$। (ख) $\cos^2(\theta/2) \times
\cos^2(\theta/2) = \cos^4(\theta/2)$। (ग) $\theta = 90^\circ$:
$(\tfrac12)^2 = \tfrac14$, अर्थात् ध्रुवक-शृंखला वाली वही संख्या। (घ)
तुच्छ रूप से $\theta = 0$; पर असली सीख यह है कि \emph{कोमल}
मध्यवर्ती माप सबसे कम कीमत लेते हैं ---
\cref{exo:b3:quantum-formalism:5} की छोटे-पग वाली सीमा, जहाँ अनेक
छोटे घूर्णन अवस्था को लगभग बिना हानि के पार करा देते हैं।
\end{solution}

\begin{solution}{exo:b3:spin-two-level:4}
(क) $2\mu_{\text{B}}B = \qty{116}{\micro eV}$: $\nu =
\qty{28}{GHz}$। (ख) $2 \times 2.79\,\mu_{\text{N}}B =
\qty{0.18}{\micro eV}$: $\nu = \qty{42.6}{MHz}$। (ग)
$k_{\text{B}}T = \qty{25.9}{meV}$ के सामने: इलेक्ट्रॉन का ध्रुवण
$\sim\num{2e-3}$, प्रोटॉन का $\sim\num{3e-6}$ --- अर्थात् तापीय
चक्रण-समुदाय लगभग पूरी तरह फेंटे हुए हैं। (घ) ESR: सूक्ष्मतरंगें;
NMR: रेडियो --- परिमाण की तीन कोटि दूर, पर एक ही भौतिकी।
\end{solution}

\begin{solution}{exo:b3:spin-two-level:5}
(क) $\dd\langle\hat S_x\rangle/\dd t = -\gamma B_0\langle\hat
S_y\rangle$, $\dd\langle\hat S_y\rangle/\dd t = \gamma
B_0\langle\hat S_x\rangle$, $\dd\langle\hat S_z\rangle/\dd t = 0$।
(ख) $\langle\hat S_x\rangle = \tfrac\hbar2\cos\omega_0t$,
$\langle\hat S_y\rangle = \tfrac\hbar2\sin\omega_0t$ ($\gamma$ के
चिह्न के अनुसार चिह्न)। (ग) दोनों (ख) से स्पष्ट हैं। (घ) अवस्था
दोनों ऊर्जा-अभिलक्षणिक अवस्थाओं का \emph{अध्यारोपण} है; उसकी आपेक्षिक
कला $\omega_0$ पर आगे बढ़ती है, और वही घूमती कला ही पुरस्सरण
\emph{है} --- स्थायी स्तर और घूमता सदिश एक ही द्वि-स्तरीय विकास के दो
पाठ हैं।
\end{solution}

\begin{solution}{exo:b3:spin-two-level:6}
(क) प्रतिस्थापन कीजिए और आधे-कोण की सर्वसमिकाएँ काम में लाइए। (ख)
$\cos\tfrac\theta2\cos\tfrac{\pi-\theta}2 +
\eu^{\iu\pi}\sin\tfrac\theta2\sin\tfrac{\pi-\theta}2 =
\cos\tfrac\theta2\sin\tfrac\theta2 -
\sin\tfrac\theta2\cos\tfrac\theta2 = 0$। (ग) विषुवत् पर,
$\varphi = 0, \pi$ ($\sigma_x$) और $\varphi = \pm\pi/2$
($\sigma_y$) पर। (घ) किसी भी विषुवतीय अक्ष के परितः $180^\circ$ का घूर्णन
--- \omterm{prop:b3:spin-two-level:rabi}{चुंबकीय अनुनाद} की $\pi$ स्पंद ही क्यूबिट का NOT फाटक है।
\end{solution}

\begin{solution}{exo:b3:spin-two-level:7}
(क) घूमते तंत्र में चालक ठहरा रहता है, जबकि तंत्र का
घूर्णन प्रभावी स्थैतिक क्षेत्र में से $\omega/\gamma$ घटा देता है
(वही हिसाब, जो किसी घूमते तंत्र के जड़त्वीय बल का होता है)। (ख) केवल
$B_1$ बचता है: ब्लॉख सदिश उस (नियत)
अनुप्रस्थ $B_1$ अक्ष के परितः $\Omega_1 = |\gamma|B_1$ पर पुरस्सरण
करता है --- अर्थात् स्थिर पलटाव। (ग) $t_\pi = \pi/\gamma B_1 = 1/(2 \times 42.58\,\text{
MHz/T} \times 10\,\unit{\micro T}) = \qty{1.2}{ms}$। (घ) एक तेज़,
सिकुड़ता शंक्वाकार सर्पिल: $z$ के परितः \qty{128}{MHz} पर पुरस्सरण,
और उसके साथ \qty{426}{Hz} पर धीरे-धीरे नीचे झुकाव --- ध्रुव से ध्रुव
तक हज़ार चक्करों का पेंच।
\end{solution}

\begin{solution}{exo:b3:spin-two-level:8}
(क) $\nu = \Delta E/h = \num{5.9e-6}/\num{4.14e-15} =
\qty{1.42e9}{Hz}$; $\lambda = c/\nu = \qty{21.1}{cm}$। (ख) $\Delta
\nu \sim 1/2\pi\tau \sim \qty{e-15}{Hz}$: बिलकुल नगण्य ---
अतः हर प्रेक्षित चौड़ाई गति है। (ग) $\Delta\nu =
\nu\,(2v/c) = 1420\,\text{MHz} \times \num{1.5e-3} \approx
\qty{2}{MHz}$: अर्थात् किसी घूमती चकती की दो-सींगों वाली रूपरेखा। (घ)
21 cm वहीं बोलती है जहाँ हाइड्रोजन परमाणविक हो (गरम, विरल); जबकि CO
वहाँ बोलती है जहाँ हाइड्रोजन युग्मित होकर अदृश्य H$_2$ बन चुका हो
(ठंडा, सघन): दोनों मिलकर किसी मंदाकिनी का पूरा अंतरतारकीय
माध्यम गिन लेते हैं।
\end{solution}

\begin{solution}{exo:b3:spin-two-level:9}
(क) दो विन्यास, सुरंगन से युग्मित: नीची ऊर्जा पर समष्टि
द्विविमीय है --- अमोनिया भेस में चक्रण-$\tfrac12$ ही \emph{है}। (ख) $\nu = \Delta E/h = \qty{24}{GHz}$, $\lambda =
\qty{1.25}{cm}$। (ग) पुंज को \emph{ऊपरी} अवस्था की अधिक जनसंख्या के
साथ आना चाहिए (जिसे स्थिरवैद्युत विक्षेपण से छाँटा जाता है) ---
अर्थात् जनसंख्या-प्रतीपन, जो तापीय रूप से मिल ही नहीं सकता, क्योंकि
बोल्ट्समान सदा नीचे वाले स्तर का पक्ष लेता है। (घ) जनसंख्या-प्रतीपन से
उद्दीपित प्रवर्धन --- अर्थात् लेज़र का कार्यकारी सिद्धांत, जो पहले
सूक्ष्मतरंग आवृत्तियों पर दिखा दिया गया।
\end{solution}

\begin{solution}{exo:b3:spin-two-level:10}
(क) $\mu_{\text{B}} \times \qty{10}{T} = \qty{0.58}{meV}$ --- अर्थात्
\cref{exo:b3:hydrogen-atom:11} का
$\alpha^2E_{\text{I}} \approx \qty{0.7}{meV}$ वाला पैमाना: सूक्ष्म
संरचना \omterm{def:b3:spin-two-level:spin}{चक्रण} का उस गति-जनित क्षेत्र से मिलना ही \emph{है}। (ख) $\Delta E = hc\,\Delta\lambda/
\lambda^2 = \qty{2.1}{meV}$: अर्थात् आंतरिक क्षेत्र $\Delta E/2\mu_{
\text{B}} \approx \qty{18}{T}$। (ग) भीतरी क्षेत्र
$Z^4$ की तरह बढ़ता है (क्षेत्र में नाभिकीय आवेश का घन, और एक बार और
कक्षा-त्रिज्या में): अतः भारी परमाणुओं की सूक्ष्म संरचना जेब के
वर्णक्रमदर्शी से भी दिख जाती है। (घ) $j = \tfrac32$: चार अवस्थाएँ; $j = \tfrac12$: दो
--- कुल छह, अर्थात् ठीक $2 \times (2l + 1)$, $l = 1$ के लिए।
\end{solution}

\begin{solution}{exo:b3:spin-two-level:11}
(क) $S_z$ की गिनतियाँ: $+\hbar, 0, 0, -\hbar$; और सममित संयोजन पर
$\hat S^2$ (सोपान-सर्वसमिका से) लगाने पर
$2\hbar^2$ ($S = 1$) मिलता है, तथा प्रतिसममित पर $0$ ($S = 0$)। (ख)
केवल एकक (और बीच वाली त्रिक-अवस्था) किसी गुणनफल के रूप में नहीं लिखी
जा सकती --- और एकक ही \emph{वह} अधिकतम उलझा हुआ
युग्म है। (ग) उत्सर्जन का अर्थ है कि त्रिक ऊपर है: अर्थात् समांतर-चक्रण
वाला विन्यास \qty{5.9}{\micro eV} अधिक ऊर्जावान है। (घ)
हर प्रेक्षक अकेले पूरी तरह यादृच्छिक परिणाम देखता है; और
सहसंबंध तभी प्रकट होता है जब दोनों सूचियाँ \emph{साथ लाई}
जाएँ --- किसी की स्थानिक सांख्यिकी नहीं बदलती, अतः कुछ भी संचरित नहीं
होता (\cref{rem:b3:quantum-formalism:nosignal} का तर्क)।
\end{solution}

\begin{solution}{exo:b3:spin-two-level:12}
(क) पहली $\pi/2$ ब्लॉख सदिश को विषुवत् पर लिटा देती है; और
$T$ के दौरान वह परमाणु के अपने $\nu_0$ पर पुरस्सरण करता है, जबकि
स्थानिक दोलित्र $\nu$ पर घूमता है: संचित कोण-अंतर
$2\pi\delta T$ है; और दूसरी $\pi/2$ उस कला को किसी
जनसंख्या में बदल देती है --- अर्थात् काल में बना एक व्यतिकरणमापी। (ख) $\Delta\delta \sim
1/2T = \qty{1}{Hz}$। (ग) $\qty{e-4}{Hz}/\qty{9.19e9}{Hz} \approx
\num{e-14}$ --- अर्थात् तीस लाख वर्षों में एक सेकंड। (घ) $10^5$ गुना
ऊँचे वाहक पर वही निरपेक्ष झालर-विभाजन आंशिक यथार्थता में
$10^5$ जीत लेता है: इसीलिए स्ट्रॉन्शियम और इटर्बियम की
प्रकाशिक घड़ियाँ $10^{-18}$ पर हैं, और सेकंड की पुनर्परिभाषा
प्रतीक्षा में है।
\end{solution}

\begin{solution}{pb:b3:spin-two-level:1}
\textbf{1.} $\nu = 42.58 \times 3.0 = \qty{127.7}{MHz}$; $h\nu =
\qty{5.3e-7}{eV}$।
\textbf{2.} $\Delta n/n = h\nu/2k_{\text{B}}T = \num{5.3e-7}/(2
\times 0.0267) \approx \num{e-5}$: अर्थात् दस भाग प्रति दस लाख।
\textbf{3.} $\num{6.6e22}\,\unit{cm^{-3}} \times \num{e-5} \approx
\num{7e17}$ संकेत देने वाले प्रोटॉन प्रति घन सेंटीमीटर --- प्रति
\omterm{def:b3:spin-two-level:spin}{चक्रण} दुर्बल, पर संख्या में प्रबल।
\textbf{4.} श्टर्न--गेरलाख ने हर परमाणु को नापा और उसे
\omterm{def:b3:quantum-formalism:observable}{अभिलक्षणिक मान} से बाँट दिया; जबकि MRI $10^{23}$ चक्रणों का
\emph{तापीय औसत} सुनती है, और बोल्ट्समान उस औसत को शून्य से
माइक्रोवोल्टों के भीतर रखता है।
\textbf{5.} ध्रुवण $\propto B_0$ और प्रेरित विद्युत्वाहक बल
$\propto \omega \propto B_0$: अतः संकेत मोटे तौर पर $\propto B_0^2$ ---
यही \qty{1.5}{T} के सामने \qty{3}{T} का, और
\qty{7}{T} की शोध-मशीनों का तर्क है।
\textbf{6.} भटकी हुई प्रवणता में कोई लौहचुंबक
$F = \mu\,\partial B/\partial z$ अनुभव करता है, जो लोहे के किलोग्रामों
तक मापित होता है: अर्थात् किसी दरवाज़े में सैकड़ों न्यूटन प्रकट हो
जाते हैं --- और यही परिरक्षण तथा जाँच-सूचियों का कारण है।
\textbf{7.} $\nu_1 = \gamma B_1/2\pi = \qty{426}{Hz}$: अतः $\pi/2$
स्पंद \qty{0.59}{ms}, और $\pi$ स्पंद \qty{1.2}{ms}।
\textbf{8.} वह अनुप्रस्थ तल में \qty{127.7}{MHz} पर पुरस्सरण करता
है; और घूमते चुंबकन का अभिवाह कुंडली से गुज़रकर (फ़ैराडे) ठीक उसी
आवृत्ति पर कोई रेडियो विद्युत्वाहक बल प्रेरित करता है ---
अर्थात् कच्चा MR संकेत।
\textbf{9.} हर ऊर्जा-विनिमयी पलटाव पलटे हुए \omterm{def:b3:spin-two-level:spin}{चक्रण} की कला भी
यादृच्छिक कर देता है: अतः जो कुछ $T_1$ बनाता है वह $T_2$ में भी
योगदान करता है, और विकलन के अपने अलग रास्ते भी हैं --- संसक्ति पहले
मरती है।
\textbf{10.} $\tau$ पर हर धावक मुड़ जाता है: जो तेज़ हैं और सबसे आगे
हैं, उन्हें अब सबसे अधिक लौटना है; और $2\tau$ पर सब एक साथ
आरंभ-रेखा पार करते हैं --- विकलन उलट जाता है और
प्रतिध्वनि बज उठती है।
\textbf{11.} केवल \emph{स्थैतिक} भाग: जिस \omterm{def:b3:spin-two-level:spin}{चक्रण} ने कोई
अचर (चाहे गलत) आवृत्ति बनाए रखी, वह अपनी कला ठीक-ठीक फिर से नाप लेता
है; जबकि उतार-चढ़ाव करते आणविक क्षेत्र दोनों अर्धांशों के बीच बदल जाते
हैं और उलटते नहीं --- उनका क्षय ही सच्चा, ऊतक-विशिष्ट $T_2$ है।
\textbf{12.} जल्दी और बार-बार नमूना लीजिए (छोटा TR, छोटा TE) तो
चर्बी का छोटा $T_1$ चमक उठता है; और देर तक रुककर देर से प्रतिध्वनि
लीजिए तो लंबे $T_2$ वाले तरल चमकते हैं: अर्थात् अनुक्रम की घड़ी की
सेटिंग ही चुनती है कि कौन-सा शिथिलन-नियतांक चित्र बनाएगा।
\textbf{13.} $\dd\nu/\dd z = 42.58\,\unit{MHz/T} \times
0.04\,\unit{T/m} = \qty{1.7}{kHz/mm}$।
\textbf{14.} वही परत, जहाँ स्थानिक लार्मर आवृत्ति उस बैंड में पड़े:
मोटाई $\qty{2}{kHz}/\qty{1.7}{kHz/mm} \approx
\qty{1.2}{mm}$ --- अर्थात् एक छाँटी हुई परत।
\textbf{15.} कोई फूरिये रूपांतर: आवृत्ति स्थान का नाम है, अतः
प्रतिध्वनि का वर्णक्रम उस परत का प्रक्षेप ही \emph{है}।
\textbf{16.} $256^2 \times 12 \approx \qty{100}{kB}$; और $T_1$ की
कोटि के हर पुनरावृत्ति-काल में एक ही संकेतित पंक्ति मिलने पर $256$
पंक्तियों में मिनट लग जाते हैं --- इसीलिए रोगियों से स्थिर रहने को
कहा जाता है।
\textbf{17.} प्रवणता-कुंडलियाँ \qty{3}{T} के भीतर किलोऐंपियर-पैमाने
की स्विच होती धाराएँ ढोती हैं: हर स्विच पर लाप्लास बल उन्हें उनके
आधारों से पीट देता है --- हर क्रमवीक्षण का मशीनगन-सा शोर।
\textbf{18.} एक्स-किरणें इलेक्ट्रॉन-घनत्व की छाया बनाती हैं --- हड्डी
बनाम हवा --- और आयनकारी मात्रा भी छोड़ जाती हैं; जबकि MRI प्रोटॉन-घनत्व
और दो शिथिलन-घड़ियाँ पढ़ती है, जो मृदु ऊतकों में भरपूर भिन्न हैं:
मस्तिष्क, उपास्थि और अर्बुद, जो एक्स-किरणों को एक-से लगते हैं,
चक्रणों के लिए अलग-अलग हैं --- और वह भी शून्य आयनकारी मात्रा पर।
\textbf{19.} हर रासायनिक परिवेश अपने प्रोटॉन को दस लाख में कुछ भाग
जितना परिरक्षित कर देता है: अतः किसी अणु के प्रोटॉन खिसकी हुई रेखाओं
की सुलझी कंघी के रूप में हाज़िरी देते हैं --- नली में ही संरचना-निर्धारण।
\textbf{20.} सक्रियता स्थानिक रक्त-प्रवाह और ऑक्सीजनन बढ़ा देती है;
अनुचुंबकीय विऑक्सीहीमोग्लोबिन पतला पड़ जाता है; स्थानिक $T_2^*$ लंबा
हो जाता है; और वह चित्रांश विचार के कुछ सेकंड बाद चमक उठता है।
\textbf{21.} \qty{84}{GHz} पर ऊतक अपारदर्शी है (मिलीमीटर
सूक्ष्मतरंगें त्वचा में मुश्किल से घुसती हैं) और इलेक्ट्रॉन का शिथिलन
माइक्रोसेकंडों का है --- जीवित शरीर में निराशाजनक, पर पदार्थों तथा
मात्रामिति में मुक्त मूलक और दोष गिनने के लिए उत्तम।
\textbf{22.} $\num{5e-7}\,\unit{eV}$, बनाम eV की बंध-ऊर्जाएँ: कुछ
भी तोड़ने के लिए एक करोड़ गुना दुर्बल --- अर्थात् अनायनकारी; जबकि
\qty{60}{keV} का एक एक्स-किरण फ़ोटॉन $10^{11}$ गुना अधिक ढोता है।
\textbf{23.} $10^{-5}$ से कोटि एक तक: अर्थात् प्रति नाभिक
$\sim10^5$ का लाभ --- और यह आवश्यक इसलिए है कि साँस में ली गई गैस जल
के प्रोटॉनों से हज़ारों गुना विरल है।
\textbf{24.} चक्रण-$\tfrac12$ कोई \emph{पूर्ण} द्वि-स्तरीय
तंत्र है (कोई रिसाव-स्तर नहीं), और वह आवेश-कोलाहल से केवल
चुंबकीय आघूर्णों के द्वारा दुर्बल रूप से युग्मित है: अर्थात् किसी
औद्योगिक पदार्थ में लंबी संसक्ति।
\textbf{25.} \qty{128}{MHz} का पुरस्सरण, $10^{-5}$ का
ध्रुवण, दो शिथिलन-घड़ियाँ और तीन प्रवणताएँ: यही चक्रण-भौतिकी किसी
जीवित मस्तिष्क का परत-दर-परत चित्र लेती है, वह भी शून्य
आयनकारी मात्रा पर --- 1922 के चाँदी के दो धब्बे, जो एक अस्पताल-विभाग
बन गए।
\end{solution}
